\chapter{Aspectos de Implementação e Avaliação}
\label{ch:implementacao}

Para que o modelo possa ser avaliado, um protótipo será construído. Esta seção descreve as tecnologias utilizadas na implementação do modelo, e a forma como será realizada esta avaliação.

%%%%%%%%%%%%%%%%%%%%%%%%%%%%%%%%%%%%%%%%%%%%%%%%%%%%%%%%%%%%

\section{Tecnologias Utilizadas}
\label{sec:tecnologias}

O protótipo do eProfile foi desenvolvido em Python \cite{python}. Python é uma linguagem genérica, orientada a objetos, de alto nível, com suporte a múltiplas arquiteturas e de código aberto.

As regras, modelos de entidade e perfis de entidade serão armazenadas no banco de dados SQLite3 \cite{sqlite3}. O SQLite3 é um banco de dados SQL auto-contido transacional, disponível em domínio público. O SQLite3 foi escolhido por apresentar bom desempenho e não necessitar de um servidor separado para sua implantação.

A comunicação entre o eProfile, o gerenciador de trilhas e as aplicações é feita através de {\it web services} \citep{fielding00}.

%%%%%%%%%%%%%%%%%%%%%%%%%%%%%%%%%%%%%%%%%%%%%%%%%%%%%%%%%%%%

\subsection{Motor de Inferência}
\label{sec:motor_inferencia}

A Figura~\ref{fig:dados_inferencia} apresenta o fluxo de dados para a inferência do perfil de uma entidade. No centro desta operação está o agente raciocinador, que utiliza um motor de inferência externo para a execução das regras. Foram estudados os seguintes raciocinadores que oferecem suporte a esta operação:

\begin{itemize}

\item {\bf Pellet:} é um racionador com suporte a OWL DL, que oferece bom desempenho e um {\em middleware} extensível. O Pellet é escrito em Java e pode ser utilizado para raciocinar com tipos de dados definidos pelo usuário e provê suporte para depuração em ontologias \cite{sirin05};

\item {\bf FaCT++:} é um raciocinador de lógica proposicional criado como uma plataforma para utilização de algoritmos de {\em tableaux} semântico e técnicas de otimização. Ele incorpora a maior parte das técnicas de otimização padrão \cite{tsarkov06};

\item {\bf HermiT:} é o primeiro raciocinador OWL que se baseia em um cálculo {\em hypertableaux}, o qual provê desempenho superior a qualquer algoritmo anterior \cite{motik09}. O HermiT é disponibilizado como um plugin do Protegé, como uma ferramenta de linha de comando ou como uma biblioteca Java.

\end{itemize}

Devido ao conjunto de funcionalidades e ao desempenho superior, o FaCT++ foi escolhido como o raciocinador utilizado pelo eProfile.

%%%%%%%%%%%%%%%%%%%%%%%%%%%%%%%%%%%%%%%%%%%%%%%%%%%%%%%%%%%%

\subsection{Linguagem de Ontologia}
\label{sec:ling_ontologia}

Para o armazenamento e processamento das ontologias, foram estudadas as linguagens de representação de conhecimento CycL \cite{lenat89}, RDF/S \cite{rdfs}, DAML+OIL \cite{daconta03} e OWL \cite{daconta03}.

% CycL
A linguagem \textbf{CycL} é a linguagem utilizada pelo projeto \emph{Cyc} \cite{lenat89}. O Cyc é um projeto de inteligência artificial que tem o objetivo de montar uma ontologia e base de conhecimento a partir do senso comum, com o propósito de permitir que aplicações de inteligência artificial realizem raciocínio semelhante ao humano. A linguagem CycL é uma linguagem declarativa baseada em lógica de primeira ordem clássica, com extensões para operadores modais e lógica de ordem superior. A base de conhecimento é dividida em micro-teorias, que são coleções de conceitos e fatos tipicamente pertencentes a um domínio particular de conhecimentos e, ao contrário da base de conhecimento como um todo, cada micro-teoria deve estar livre de contradições.

% RDF/S
\textbf{RDF} (\emph{Resource Description Framework}) é uma linguagem utilizada para representação de metadados a respeito de recursos na Web \cite{rdf}. No entanto, o conceito de ``recurso da Web'' pode ser generalizado, de forma que o RDF pode ser utilizado para representar informações a respeito de coisas que podem ser identificadas na Web, mesmo que elas não possam ser carregadas da Web. Recursos podem ser recursos eletrônicos, como arquivos, ou conceitos, como pessoas. O RDF é complementado pelo \textbf{RDF/S} (\emph{RDF/Schema}), uma extensão semântica do RDF. O RDF/S é uma linguagem de descrição de vocabulário que provê mecanismos para descrever grupos de recursos relacionados e os relacionamentos entre estes recursos \cite{rdfs}.

% DAML+OIL
O \textbf{DAML+OIL} é uma linguagem de representação de conhecimento baseada nas linguagens DAML e OIL. O DAML (\emph{DARPA Agent Markup Language}) foi desenvolvido pelo Centro de Pesquisas Avançadas da Defesa Americana e é uma linguagem de marcação de agentes baseada no RDF. Já o OIL (\emph{Ontology Inference Layer}) é uma linguagem com os mesmos objetivos, criada por um grupo de universidades europeias. Ambos os esforços foram unidos na criação do DAML+OIL, que por sua vez deu origem ao OWL \cite{daconta03}.

% OWL
O \textbf{OWL} é a linguagem de ontologia mais expressiva definida ou em definição para a Web Semântica \cite{daconta03}. O OWL utiliza uma URI (\emph{Uniform Resource Identifier}) para os nomes, e a estrutura de descrição disponibilizada pelo RDF. O OWL está uma camada acima do RDF e do RDF/S, e adiciona mais vocabulário para descrever propriedades e classes, tais como relações entre classes, cardinalidade, igualdade, tipagem mais avançada de propriedades, características de propriedades e classes enumeradas \cite{owl-guide}. O OWL está disponível em três sub-linguagens: \emph{OWL Lite}, \emph{OWL DL} e \emph{OWL Full}, cada uma aumentando o nível de complexidade e expressividade. Um grande número de ferramentas com capacidade para edição e raciocínio semântico com suporte para OWL está disponível no mercado.

% Escolhido - OWL
Devido à expressividade e ao suporte de ferramentas externas, o OWL DL foi escolhido como a linguagem de ontologias utilizada pelo eProfile para modelagem das ontologia de perfis e de trilhas.

%%%%%%%%%%%%%%%%%%%%%%%%%%%%%%%%%%%%%%%%%%%%%%%%%%%%%%%%%%%%

%\section{Metodologia de Avaliação}
%\label{sec:avaliacao}

%O trabalho apresenta um modelo que auxilia aplicações a montarem perfis de entidades utilizando inferência em trilhas. O eProfile é uma proposta de infraestrutura e, portanto, sua avaliação é difícil, uma vez que não é visível para o usuário final. Segundo \citetexto{edwards03}, somente é possível avaliar as funcionalidades de uma proposta de infraestrutura construindo aplicações que a utilizem e então avaliá-las, de modo a obter uma avaliação indireta do {\em framework}.

%Assim sendo, para a avaliação desta proposta será realizada uma abordagem sob duas perspectivas: {\bf funcionalidade} e {\bf cenário}. Na primeira, a avaliação se dau sobre as funcionalidades propostas, isto é, se o protótipo é funcional e atende o modelo proposto no Capítulo~\ref{ch:modelo}. Na segunda, foi desenvolvido um protótipo de aplicação, que seguiu um roteiro definidos em um cenário. Foi, então, verificado se o eProfile conseguiu cumprir as premissas descritas no cenário e se a utilização do eProfile trouxe aperfeiçoamento das informações geradas no cenário descrito.

%Para o teste de funcionalidades, foram verificadas as premissas descritas abaixo.

%\begin{itemize}

%\item o eProfile é capaz de se conectar a uma aplicação e um gerenciador de trilhas;
%\item a aplicação pode registrar uma nova entidade no eProfile;
%\item a aplicação pode registrar um novo modelo de entidade no eProfile;
%\item a aplicação pode registrar novas regras no eProfile;
%\item o eProfile busca a trilha da entidade no gerenciador de trilhas;
%\item o eProfile traduz a trilha da entidade do formato do gerenciador de trilhas para seu formato próprio;
%\item o eProfile é capaz de utilizar as regras para realizar a inferência nas trilhas de uma entidade, e montar um perfil de acordo com o modelo de entidade;
%\item o eProfile é capaz de realizar todos os testes acima para múltiplos gerenciadores de trilhas, múltiplas aplicações e múltiplas entidades.

%\end{itemize}

%Para a realização do teste de cenário, foi desenvolvido um cenários onde foram simuladas atividades realizadas por múltiplas aplicações e múltiplas entidadas. Através da análise das informações geradas pelas aplicações simuladas no cenário, foi possível verificar que o eProfile potencializou a capacidade de personalização das aplicações. O cenário está descrito no 
